\section{T3真题刷题阶梯}
按下面四档推进。每题先写对象表、规则表和输出格式，再独立实现；完成后记录最早出错的位置，以及能暴露该错误的最小输入。表中的主要坑点可作为自测方向。

\subsection{第一档：纯规则模拟}
目标是完整读入题面信息，准确执行每条规则，并完成全部输出。
\begin{longtable}{P{5.8cm}P{3cm}P{4.9cm}}
\toprule 真题入口 & 对应构件 & 主要坑点\\\midrule\endhead
\csp{CSP201403C}{命令行选项} & 分词与状态机 & 带参选项、遇到非法项时的停止位置、重复选项。\\
\csp{CSP201604C}{路径解析} & 目录栈 & 空路径、连续斜杠、根目录上移、每条查询独立。\\
\csp{CSP201709C}{JSON查询} & 扫描与嵌套对象 & 字符串转义、点号路径、对象与字符串的输出类型。\\
\bottomrule\end{longtable}
路径解析的空行输入可专门用getline处理；基础构件中的路径题约定非空，做真题时将空路径按当前目录处理。JSON查询先完整解析对象，再按查询路径逐段查找。

\subsection{第二档：规则与层次结构}
目标是把文本层次保存为明确的状态，区分整行规则和行内规则。
\begin{longtable}{P{5.8cm}P{3cm}P{4.9cm}}
\toprule 真题入口 & 对应构件 & 主要坑点\\\midrule\endhead
\csp{CSP201703C}{Markdown} & 行状态与行内扫描 & 空行结束区块、标题与列表、链接和强调的嵌套。\\
\csp{CSP201809C}{元素选择器} & 树与祖先链 & 标签和ID的大小写规则、多级严格祖先。\\
\csp{CSP201509C}{模板生成系统} & 字段提取与替换 & 标记空格、缺失变量、替换值按原题规则输出。\\
\bottomrule\end{longtable}
模板生成系统按原题进行变量替换；前面的教学构件另练递归模板依赖与三色判环。做题前先将两种模型各自的输入和展开规则列清楚。

\subsection{第三档：复杂逻辑与语法解析}
目标是明确递归边界、对象生命周期和多表之间的关系。
\begin{longtable}{P{5.8cm}P{3cm}P{4.9cm}}
\toprule 真题入口 & 对应构件 & 主要坑点\\\midrule\endhead
\csp{CSP201912C}{化学方程式} & 递归计数 & 前置系数、括号后数量、多位数与元素名。\\
\csp{CSP202104C}{DHCP服务器} & 状态机与时间轴 & 租约到期、报文接收对象、地址状态更新顺序。\\
\csp{CSP202206C}{角色授权} & 集合与倒排索引 & 同一角色满足三项条件、通配规则、当次用户组。\\
\bottomrule\end{longtable}
化学方程式可按“元素、数字、单层括号、嵌套括号”递进；DHCP服务器按每种报文画状态转移表；角色授权先用小角色集合手算并集和交集，再实现完整查询。

\subsection{第四档：近年综合模拟与性能}
目标是保持规则完整，同时控制扫描、查找与复制的成本。
\begin{longtable}{P{5.8cm}P{3cm}P{4.9cm}}
\toprule 真题入口 & 对应构件 & 主要坑点\\\midrule\endhead
\csp{CSP202305C}{解压缩} & 字节流与扫描指针 & 变长长度、小端序、重叠回溯、每行8字节。\\
\csp{CSP202309C}{梯度求解} & 表达式与递推 & 运算栈次序、乘法求导、取模与负数。\\
\csp{CSP202312C}{树上搜索} & 树与动态候选集 & 每轮权重变化、并列编号、查询状态重置。\\
\csp{CSP202412C}{缓存模拟} & 映射与淘汰状态 & 组号计算、访问次序、替换与写回规则。\\
\bottomrule\end{longtable}
完成本档后，挑一题重新限时独立实现。复盘时保留三项记录：函数职责划分、最早出现的错误、满数据时最耗时的操作。下一次练习围绕其中一项改进，再用同一组边界数据检查。
